\documentclass[11pt,a4paper]{article}
\usepackage[margin=2.3cm]{geometry}
\usepackage{amsmath,amssymb}
\usepackage{booktabs}
\usepackage{enumitem}
\usepackage{xcolor}
\usepackage{listings}
\usepackage[colorlinks=true,linkcolor=blue!50!black,urlcolor=blue!50!black]{hyperref}
\usepackage{titlesec}
\titleformat{\section}{\Large\bfseries\color{black!80}}{}{0pt}{}
\titleformat{\subsection}{\large\bfseries\color{black!70}}{}{0pt}{}
\lstset{basicstyle=\ttfamily\small,frame=single,framerule=0.3pt,columns=fullflexible,
        breaklines=true,escapeinside={(*}{*)}}

\title{\vspace{-1.2cm}\textbf{LETMA 2026 Summer School}\\[2pt]
\textbf{Machine Learning on Geophysical Data --- Labs: Solutions \& Expected Results}\\[6pt]
\large Instructor handoff: the answers to the exercises and the numbers every run reproduces}
\author{\small LETMA 2026 Summer School}
\date{}

\begin{document}
\maketitle
\vspace{-6mm}

This is the companion of the student handout: for each lab it gives the exercise solutions
and the expected results. Every number below was produced by executing the
\texttt{*\_solution} notebooks end to end on a MIG slice (\texttt{1g.20gb}) of an H100 on the
Hedy cluster; the executed notebooks contain every figure and animation, so any of these
results can be shown live without re-running a cell.

\section{Lab 1 --- FNO for the 1-D Burgers equation}

\subsection*{Exercise F1 (spectral convolution)}
\begin{lstlisting}
def forward(self, x):
    n = x.shape[-1]
    xh = torch.fft.rfft(x, dim=1, norm='ortho')       # (B, C, N//2+1)
    m = min(self.modes, xh.shape[-1])                 # truncation (H7)
    out = torch.einsum('bik,iok->bok', xh[..., :m], self.weight[..., :m])
    return torch.fft.irfft(out, n=n, norm='ortho')    # back to physical space
\end{lstlisting}
\textbf{Answer to the question.} Keeping $m=16$ modes loses almost nothing because the target
fields' energy is concentrated in the first Fourier modes (H7, checked on data): mode $m=16$
already reconstructs a target to $\approx 0.3\%$ relative error.

\subsection*{Exercise F2 (forward pass)}
Lift $a(x)\oplus x$ to $w=32$ channels ($P$), then four blocks
$v_{\ell+1}=\mathrm{GELU}(K_\ell v_\ell + W_\ell v_\ell)$, then $Q$ (width $128$) back to one
channel. \textbf{Answer:} the computation happens in the learned 32-channel latent space
--- the blocks never see $u$ again after the lift.

\subsection*{Expected results}
\begin{center}
\begin{tabular}{@{}ll@{}}
\toprule
validation rel-$L^2$ (60 epochs, AdamW+cosine) & $\approx 0.083$\\
test rel-$L^2$ (untouched split) & mean $0.076$, median $0.075$, worst $0.179$\\
zero-shot super-resolution $N{=}64\to N{=}256$ & $0.073$ (no retraining)\\
optimizer race to $L^*$ & Adam \textbf{79 steps}, SGD+momentum \textbf{604 steps},\\
 & plain SGD still far above $L^*$ after 6000 steps\\
two-seed minima & two distinct minima, three Adam runs fall in separate basins\\
corner plot (best of 20 configs, 25 epochs) & modes 16, width 48, 6 layers, proj 128 $\to$ $0.090$\\
\bottomrule
\end{tabular}
\end{center}

\section{Lab 2 --- FFNO for infrasound boundary-layer propagation}

\subsection*{Exercise F2 (signal spectrum)}
\begin{lstlisting}
def signal_spectrum(trace, dt):
    f = np.fft.rfftfreq(len(trace), d=dt)     # frequencies of the rfft bins
    return f, np.fft.rfft(trace)              # complex spectrum
\end{lstlisting}
\textbf{Answer to the question.} The energy of a recorded trace lives inside the band
$[0.3, 4]$ Hz propagated by the solver (H4); about 64 of the 129 time Fourier modes carry the
whole band --- exactly the truncation the factorized operator makes (H13).

\subsection*{Expected results}
\begin{center}
\begin{tabular}{@{}ll@{}}
\toprule
demo atmosphere & rays return to the ground ($\approx$19 arrivals within 11 km);\\
 & hodograph arrivals at $t = r/c_\mathrm{ref} + $ wavelet delay\\
training (240 epochs, 200 atmospheres) & validation rel-$L^2 \approx 0.30$\\
test (28 unseen atmospheres) & mean $0.266$, worst $0.597$\\
loss family & rel-$L^2$ better than MSE on validation (quiet atmospheres served)\\
overfitting demo (8 atmospheres) & train $0.56$ vs validation $1.31$ (gap $\times 2$)\\
optimizer race & Adam reaches $L^*$ in $\approx 250$ steps; fixed-step methods\\
 & remain near $1.0$ (near-silence initialization) after 4000 steps\\
\bottomrule
\end{tabular}
\end{center}

\section{Lab 3 --- GNN for a sensor network}

\subsection*{Exercise G1 (message passing)}
\begin{lstlisting}
def message(self, h_i, h_j, pos_i, pos_j):
    rel = pos_j - pos_i
    return self.msg(torch.cat([h_i, h_j, rel, rel.norm(dim=-1, keepdim=True)], -1))
def forward(self, h, pos, edge_index):
    return h + self.upd(torch.cat([h, self.propagate(edge_index, h=h, pos=pos)], -1))
\end{lstlisting}
\subsection*{Exercise G2 (the locator)}
\begin{lstlisting}
h = self.encoder(data.x)                       # one 64-dim state per station
for layer in self.layers:                     # 4 message-passing layers
    h = layer(h, data.pos, data.edge_index)
pooled = global_mean_pool(h, data.batch)      # invariant event code
return self.decoder(pooled)                   # normalized hypocentre
\end{lstlisting}

\subsection*{Expected results}
\begin{center}
\begin{tabular}{@{}ll@{}}
\toprule
smoothing-operator spectrum (demo graph) & $\rho = 0.844$: the slowest non-mean mode keeps\\
 & $0.51$ of its amplitude at depth 4, $0.066$ at depth 16\\
depth experiment (validation error) & 1 layer $0.47$ \; 2: $0.43$ \; 4: $0.43$ \; 8: $0.37$ \; 16: $0.39$\\
locator (60 epochs) & validation $\approx 0.31$ (normalized; $1.0 \sim 20$ km)\\
epicentral error (50 test events) & median $4.87$ km, mean $5.08$, 90\% $8.23$\\
graph-free MLP baseline & median $6.92$ km (pays for every missing station)\\
\bottomrule
\end{tabular}
\end{center}
\textbf{Answers to the questions.} The aggregation must be a mean/sum because only those
treat neighbours unordered (permutation invariance). The GNN survives station dropout because
the graph and the masked pooling adapt to any node count; the fixed-size MLP has one slot per
station and loses the information of every station that is down. Node states become
interchangeable when the Dirichlet energy collapses: one scalar per depth is the diagnostic.

\section{Optional extensions}
\begin{itemize}[leftmargin=*,itemsep=2pt]
\item \textbf{Darcy scale-up:} a 123k-parameter FFNO2d, 200 epochs on 2304 pairs at $96\times96$:
test rel-$L^2$ $0.0083$ (MIG, $\approx 9$ min).
\item \textbf{Labs 1 and 3 in JAX} (Equinox + Optax): FNO val $0.0969$, test $0.0867$,
zero-shot at $N{=}256$: $0.0889$; GNN epicentral median $3.05$ km vs MLP $7.04$ km. Same
physics, different framework --- the notebooks need neither torch nor PyG (data shipped as
arrays).
\end{itemize}

\section{How to re-run everything}

\subsection*{Connect and work on Hedy (one SSH hop)}
\begin{lstlisting}
ssh -Y hedy                       # the single hop (-Y: X11 forwarding)

# either use the pre-built shared environment (nothing to install):
source /home/projects/m26266/letma2026-ml-labs/env/fno-gnn/bin/activate

# or install the requirements into your own venv:
python3 -m venv ~/envs/fno-gnn && source ~/envs/fno-gnn/bin/activate
pip install -r /home/projects/m26266/letma2026-ml-labs/requirements.txt

# copy the labs into YOUR home folder on Hedy (env/ excluded: the venvs are 5+ GB):
rsync -a --exclude env /home/projects/m26266/letma2026-ml-labs/ ~/letma2026-ml-labs/

# GPU session on a MIG slice: ONE command allocates the slice, forwards X11
# and drops you into an interactive shell ON the GPU node:
salloc \
    --account=m26266 \
    --partition=gpu \
    --reservation=test_mig_letma \
    --nodelist=m-gpu[07-08] \
    --gres=gpu:nvidia_h100_80gb_hbm3_1g.20gb:1 \
    --nodes=1 \
    --ntasks=1 \
    --cpus-per-task=8 \
    --time=01:00:00 \
    bash -c '
        NODE=$(scontrol show hostnames "$SLURM_JOB_NODELIST" | head -n 1)
        echo "Allocated job: $SLURM_JOB_ID on $NODE"
        exec ssh -Y -C -t "$NODE" \
            "srun --jobid=$SLURM_JOB_ID --overlap --ntasks=1 --pty bash -i"
    '
# release afterwards with: scancel <jobid>

# on the GPU node, first CHECK YOUR SLICE - one MIG device with ~20 GB:
nvidia-smi
#   MIG devices:
#   GPU  GI  CI  MIG |              Shared Memory-Usage
#     0   3   0   0  |               9MiB / 20096MiB     <- your 1g.20gb slice
# (if the MIG section is absent or shows another size, exit and resubmit the salloc)

# open a notebook with Euporie, the notebook TUI (Hedy has no SSH port
# forwarding, so there is no browser interface on the cluster; jupyter lab
# remains the normal interface on the participants' own machines):
cd ~/letma2026-ml-labs/notebooks/labs
euporie-notebook --kernel-name python3 lab_fno_solution.ipynb
\end{lstlisting}
Launch time: about 5 seconds on a warmed node, 1--2 minutes on a cold one.
Everything (euporie, jupyter, torch, numpy, matplotlib, ...) is paged in from the shared
filesystem on the first launch of a node; the screen shows \texttt{Loading...} /
\texttt{No Kernel} meanwhile --- this is not a freeze. \textbf{Before the session, warm up
both nodes} (the Linux page cache is shared by all users of a node, so one run makes every
participant's launch quick):

\begin{lstlisting}
ssh hedy   # no -Y needed here: this submits a batch job, nothing is displayed
sbatch /home/projects/m26266/letma2026-ml-labs/slurm/prewarm_nodes.sbatch   # ~3 min, both nodes
\end{lstlisting}

Participants then wait only for the kernel badge \texttt{Python 3 (ipykernel)} to appear
in the bottom-right status bar. The \texttt{Rendering... (61\%)} indicator may stay at a
partial value: it counts the whole notebook, and cells below the visible area are rendered
lazily when you scroll.

Euporie quick reference (verified against the installed version): \texttt{Up/Down} (or
\texttt{j/k}) select cells, \texttt{Home/End} first/last, \texttt{PgUp/PgDn} $\pm$5;
\texttt{Enter} enters edit mode, \texttt{Esc} leaves it; \texttt{Ctrl+Enter} runs the selected
cell, \texttt{Shift+Enter} runs and selects the next, \texttt{Alt+Enter} runs and inserts below;
\texttt{a}/\texttt{b} insert cells above/below, \texttt{d d} deletes, \texttt{z} undeletes,
\texttt{x}/\texttt{c}/\texttt{v} cut/copy/paste, \texttt{m}/\texttt{y} set the cell type;
\texttt{i i} interrupts the kernel, \texttt{0 0} restarts it; \texttt{Ctrl+S} saves,
\texttt{Ctrl+O} opens, \texttt{Ctrl+Q} quits (prompting to save unsaved changes); \texttt{Tab}
cycles the focus to the command bar for the rarer actions. Figures render inline (terminal
graphics on Kitty/Sixel terminals, character previews elsewhere). The \texttt{python3} kernel
is the activated venv's interpreter --- the validated one. The JAX extensions use the sibling
venv \texttt{env/jax-hedy} (activate it the same way). To work on a GPU: \texttt{ssh -Y hedy},
then run the single \texttt{salloc} command of the listing below: it allocates the
MIG slice, forwards X11 with \texttt{ssh -Y} into your allocated node, and drops you
into an interactive shell there via \texttt{srun --pty} (which is why
\texttt{srun --x11} is not needed here).

\subsection*{Headless re-validation of the whole set}
On Hedy, from the shared folder, one command re-executes all six lab notebooks plus the
extensions headlessly (see \texttt{slurm/PROCEDURE\_TESTS.md}):
\begin{lstlisting}
sbatch slurm/test_all_mig.sbatch
\end{lstlisting}
Mandatory MIG directives are already in the script. Total measured time: labs
$\approx 5$ min, Darcy $\approx 9$ min, JAX $\approx 7$ min.

\end{document}
